\textbf{Ferity} jsou sloučeniny oxidu železa s oxidy dalších kovů. Mají malou elektrickou vodivost, a tím malé ztráty vířivými proudy; jsou tvrdé, křehké, chemicky netečné a nenavlhají.
\begin{itemize}
\item \textbf{Spinelové ferity} $\mathrm{MO\cdot Fe_2O_3}$: magnetit $\mathrm{FeO\cdot Fe_2O_3}$, nikelnatý a kobaltnatý ferit. Jejich vlastnosti závisejí na obsazení podmřížek; normální struktury (např. Zn a Cd ferit) nemají ferimagnetické uspořádání, inverzní je mají.
\item \textbf{Smíšené ferity Mn--Zn, Ni--Zn, Cu--Zn:} kombinují oxidy více kovů; díky lepším magnetickým vlastnostem patří k nejpoužívanějším.
\item \textbf{Ferity typu magnetoplumbitu:} šesterečná struktura, např. $\mathrm{PbO\cdot6Fe_2O_3}$; také ferity s Ba či Sr.
\item \textbf{Granátové ferity:} obsahují prvky jako Y, Gd nebo Sc, mají granátovou strukturu se třemi podmřížkami.
\end{itemize}
